\section{Materiais e métodos}
\label{sec:material_metodos}

\subsection{Materiais}

\subsubsection{Nariz eletrônico}
\label{subsec:sistema_nariz}

O nariz eletrônico foi desenvolvido para registrar a assinatura de compostos
voláteis.
Foi composto por microcontrolador ESP32, seis sensores de gases MQ, sensor
capacitivo de umidade do solo (Soil Moisture), sensor BMP280 de temperatura e
pressão, bomba de vácuo, válvulas, tubo de PVC acoplado ao vaso, mangueiras,
câmara sensorial de vidro
e estrutura externa de MDF. O ESP32 registrou os sinais e controlou a bomba
e as válvulas; uma interface web acessada por Wi-Fi permitiu iniciar e
acompanhar as aquisições.
O BMP280 forneceu temperatura e pressão medidas no equipamento, não uma
medida de umidade relativa do ar na câmara sensorial. O sensor capacitivo
Soil Moisture forneceu a leitura usada como índice operacional do substrato.
Esse índice não equivale a uma medida calibrada de água no solo.

\begin{figure}[H]
    \centering
    \includegraphics[width=0.30\textwidth]{figures/completo.jpeg}
    \caption{Estrutura externa do nariz eletrônico.}
    \label{fig:estrutura_nariz_eletronico}
\end{figure}

\paragraph{Matriz de sensores de gases}
\label{subsubsec:matriz_sensores}

A matriz foi constituída por MQ-2, MQ-3, MQ-7, MQ-8, MQ-135 e MQ-138,
posicionados na câmara de vidro e expostos à mesma amostra. Esses sensores
semicondutores de óxido metálico apresentam sensibilidades
sobrepostas. Por isso, a classificação utilizou o padrão conjunto das seis
respostas, sem atribuir a um sensor isolado a identificação química de um
composto. As curvas de sensibilidade dos fabricantes foram empregadas apenas
para descrever as respostas nominais dos modelos MQ. As leituras registradas
foram sinais instrumentais, sem conversão para concentrações de compostos.

\begin{figure}[H]
    \centering
    \includegraphics[width=0.30\textwidth]{figures/matriz.png}
    \caption{Matriz dos seis sensores de gases na câmara sensorial.}
    \label{fig:matriz_sensores}
\end{figure}


\subsubsection{Material vegetal e inóculo}
\label{subsubsec:material_vegetal}

Foram utilizados dez vasos com plantas de soja e juvenis de segundo
estádio (J2) de \textit{Heterodera glycines} para a inoculação.
Os vasos constituíram as unidades biológicas do experimento.

\subsection{Metodologia}

\subsubsection{Metodologia experimental}
\label{subsubsec:metodologia_experimental}

\paragraph{Delineamento experimental}
\label{subsec:delineamento}

O experimento foi conduzido na Estação de Pesquisas Agronema, em Grandes
Rios, Paraná, Brasil, nas coordenadas $24{,}186^{\circ}$ S e
$51{,}399^{\circ}$ O. Os dez vasos foram divididos em dois grupos de cinco.
Leituras sucessivas e coletas repetidas do mesmo vaso não foram tratadas
como réplicas biológicas independentes.


\paragraph{Grupos experimentais e inoculação}
\label{subsec:classes_experimentais}

Conforme o delineamento experimental, os vasos 1 a 5 formaram o
grupo controle, considerado saudável e não inoculado. Os vasos
6 a 10 formaram o grupo inoculado e receberam cerca de
100 mil juvenis de segundo estádio (J2) de
\textit{H. glycines} por vaso.

A Figura \ref{fig:vasos_experimento} apresenta os vasos
utilizados no experimento, divididos entre os grupos
controle e inoculado.

\begin{figure}[H]
    \centering
    \includegraphics[width=0.55\textwidth]{figures/vasos_experimento.png}
    \caption{Vasos de soja utilizados no experimento,
    divididos entre os grupos controle não inoculado
    e inoculado com \textit{H. glycines}.}
    \label{fig:vasos_experimento}

    \vspace{0.1cm}
    {\footnotesize Fonte: Elaborado pelos autores (2026).}
\end{figure}

Na base de classificação, a \textbf{classe 0} identifica o grupo
inoculado no experimento e a \textbf{classe 1} identifica o controle
não inoculado. Esses rótulos representam a alocação experimental, não a
confirmação nematológica individual de cada vaso. Neste artigo, o recall
da classe inoculada corresponde ao recall da classe 0.


\paragraph{Calendário de aquisição}

As aquisições começaram cerca de 35 dias após a inoculação, em
11 de maio de 2026. A base de coletas a vácuo reuniu 36 coletas nos dias
11, 12, 13, 18 e 19 de maio de 2026, cerca de 35, 36, 37, 42 e 43 dias
após a inoculação. Dessas coletas, 17 pertenciam ao grupo inoculado e 19 ao
controle. Cada vaso foi amostrado em três ou quatro ocasiões. Houve irrigação
entre as coletas de 13 e 18 de maio, com mudança da faixa operacional do
índice de solo entre esses períodos.

\paragraph{Protocolo de coleta a vácuo}
\label{subsec:protocolo_aquisicao}

O tubo de PVC foi acoplado ao vaso para delimitar o espaço de amostragem
(Figura~\ref{fig:vasos_experimentos}). Antes da entrada da amostra, a
bomba, controlada pelo ESP32, reduzia a pressão interna do sistema para
cerca de 70~kPa, enquanto a pressão externa estava próxima de
94~kPa. O valor de 70~kPa era a referência operacional; o ESP32 comandava
o encerramento da depressurização e a abertura da válvula. Na entrada inicial da
amostra, a diferença de pressão conduzia o ar do \textit{headspace} do vaso pela
tubulação até a câmara dos sensores. Após essa admissão, o ar permanecia
retido na câmara durante a aquisição; na etapa de desaturação, o ar era
liberado. O protocolo não pressupõe fluxo contínuo de amostra durante
todo o intervalo de leitura. Uma execução
completa do ciclo, registrada sob um identificador próprio de vaso e dia, foi
definida como uma coleta. Em cada leitura, o ESP32 registrou o tempo decorrido,
os seis sinais MQ, a leitura do sensor de solo, a temperatura e a pressão.
A base analisada reuniu as coletas do protocolo de amostragem a vácuo.

\begin{figure}[H]
    \centering
    \includegraphics[width=0.55\textwidth]{figures/vasos_experimentos.png}
    \caption{Protótipo posicionado sobre o vaso de soja durante a coleta a vácuo.}
    \label{fig:vasos_experimentos}
    \vspace{0.1cm}
    {\footnotesize Fonte: Elaborado pelos autores (2026).}
\end{figure}

A Figura~\ref{fig:ciclo_coleta} apresenta as etapas de operação do sistema
com e sem pré-depressurização. O termo ``coletas a vácuo'' designa aqui as
aquisições em que a câmara foi depressurizada pela bomba e
recebeu a amostra após a abertura da válvula. Coletas de outros protocolos
não foram combinadas com elas nesta análise. A pressão registrada como
entrada dos classificadores corresponde às leituras retidas após o recorte
temporal e o filtro descritos adiante, não ao valor-alvo de 70~kPa usado
na preparação da câmara.

\begin{figure}[H]
    \centering
    \includegraphics[width=\textwidth]{figures/fluxo_en.png}
    \caption{Fluxograma dos ciclos de coleta com pré-depressurização
    e esquema do sistema de amostragem.}
    \label{fig:ciclo_coleta}
    \vspace{0.1cm}
    {\footnotesize Fonte: Elaborado pelos autores (2026).}
\end{figure}

O ciclo de referência, ilustrado na Figura~\ref{fig:ciclo_coleta}, compreendeu
quatro etapas: aquecimento dos sensores por 180~s; extração do ar até a pressão
interna atingir 70~kPa; admissão da amostra e leitura dos sensores por 600~s;
e desaturação dos sensores por 180~s. O ESP32 controlou a bomba e as válvulas.
A frequência de amostragem informada para a etapa de leitura foi de três a
quatro leituras por segundo. As etapas com tempo fixo somam 960~s; a extração
acrescenta um intervalo de duração variável. Esses tempos descrevem o ciclo de
referência, não a duração medida de cada registro.

\subsubsection{Metodologia de análise dos dados}
\label{subsubsec:metodologia_analise}

\paragraph{Pré-processamento dos sinais e janelamento}
\label{subsec:preprocessamento}

Cada planilha da base original representou uma coleta. As leituras receberam
identificadores de coleta, dia, vaso e classe. Na consolidação, foram removidos
15\% dos registros iniciais e 15\% dos registros finais de cada série. O
recorte reteve a faixa central de 70\% dos registros de cada coleta. Esse
critério usa a posição das leituras na série e não delimita, por tempo real,
a etapa nominal de 600~s. Não havia marca de abertura da válvula verificada
em cada série. As colunas usadas na análise foram convertidas para números;
registros sem valores válidos nessas colunas foram descartados.

A condição de solo seco ou úmido foi definida pelo índice operacional
\texttt{Soil\_indice\_0\_1}, obtido pela normalização mínimo--máximo da leitura
do sensor de umidade do solo na base consolidada, conforme a
Equação~\eqref{eq:indice_solo}:
\begin{equation}
    I_{\mathrm{solo}} =
    \frac{\mathrm{Soil}-\min(\mathrm{Soil})}
         {\max(\mathrm{Soil})-\min(\mathrm{Soil})},
    \label{eq:indice_solo}
\end{equation}
em que $\mathrm{Soil}$ é a leitura bruta do sensor de umidade do solo, e
$\min(\mathrm{Soil})$ e $\max(\mathrm{Soil})$ são os valores
mínimo e máximo dessa leitura na base consolidada, nessa ordem.
Esses extremos foram determinados antes das divisões de treino e teste;
portanto, as leituras destinadas ao teste participaram da
definição da escala do índice. A escala não foi reajustada dentro
de cada rodada de validação.

Leituras com $I_{\mathrm{solo}}\leq 0{,}4$ foram denominadas \emph{solo seco};
leituras com $I_{\mathrm{solo}}>0{,}4$, \emph{solo úmido}. No conjunto de
teste da partição histórica, todas as leituras de cada coleta ficaram do mesmo lado desse limiar,
permitindo classificar também a coleta inteira como seca ou úmida. O índice não
representa porcentagem calibrada de água no solo.

Para remover transientes de pressão, calculou-se
$\Delta P_t=|P_t-P_{t-1}|$ dentro de cada coleta. Quando
$\Delta P_t\geq 0{,}10$~kPa, foram excluídas a leitura do evento e até
30 leituras anteriores e 30 posteriores. Na versão estrita do filtro,
também foram removidas leituras fora da faixa definida pela mediana global
da pressão $\pm 0{,}50$~kPa, calculada após o filtro de saltos na base
consolidada. Essa mediana foi 93,69~kPa, resultando na faixa de
93,19 a 94,19~kPa. O filtro reteve leituras próximas à
pressão externa após a entrada inicial da amostra, não as leituras da
depressurização até cerca de 70~kPa. Após a
filtragem estrita, a base de coletas a vácuo continha 63.492 leituras de
36 coletas. Não foi aplicada média móvel nem outra suavização temporal
aos sinais usados pelos classificadores.
A exclusão em janelas de até 30 leituras de cada lado do evento retirou
2.276 de 32.226 leituras da classe 0 (7,06\%) e 2.945 de 39.076 da
classe 1 (7,54\%). A restrição adicional pela faixa de pressão removeu
mais 906 leituras da classe 0 (2,81\% das 32.226 originais) e 1.683 da
classe 1 (4,31\% das 39.076 originais). Considerando as duas
etapas, foram retirados 9,87\% e 11,84\% das leituras originais dessas
classes na mesma ordem.
A mediana de pressão também foi calculada antes das divisões de treino e
teste. A faixa central proporcional não corresponde de forma exclusiva
à fase nominal de exposição ao \textit{headspace}.

Na análise histórica com divisão por coleta, as entradas foram seis
sinais MQ com correção estatística, o índice de solo, a temperatura em $^{\circ}$C e a
pressão em kPa. Para cada MQ, ajustou-se com as coletas de
treinamento um \texttt{HuberRegressor} que estimou a componente associada
ao índice de solo, à temperatura e à pressão.
Os
``MQ corrigidos'' são, portanto, sinais com compensação \emph{estatística}
da associação linear estimada com essas três variáveis. Eles não resultam
de calibração física para concentração de gases ou umidade relativa do ar.
Essa compensação pertence à análise histórica por coleta; as análises LOSO
descritas adiante utilizaram sinais MQ brutos. As seis entradas MQ, o
índice de solo, a temperatura e a pressão totalizaram nove canais por
leitura na análise histórica.

Para a TCN, as leituras foram ordenadas por coleta e tempo e segmentadas
em janelas de 128 leituras, com passo de 64 leituras e sobreposição nominal
de 50\%, sem atravessar limites entre coletas. Quando necessário, uma
janela adicional foi alinhada ao fim da coleta para incluir as últimas
leituras; sua sobreposição podia superar 50\%. Cada janela teve formato de
$9\times128$ (canais $\times$ leituras) na análise histórica. As entradas
da rede foram padronizadas com parâmetros estimados nas coletas de
treinamento. Na base processada, o intervalo mediano observado entre
leituras consecutivas da mesma coleta foi de cerca de 0,24~s;
como há lacunas de aquisição, 128 leituras não equivalem sempre à mesma
duração em segundos.


\paragraph{Modelos de aprendizado de máquina}
\label{subsec:modelos}

Foram utilizados três classificadores principais:
\emph{Extra Trees} (ET), \emph{Random Forest} (RF) e uma rede
convolucional temporal (TCN). A validação principal por vaso empregou
ET; a análise histórica com separação por coleta comparou ET e TCN; e
a análise exploratória após a exclusão de cinco coletas comparou os
três modelos. A regressão logística serviu como controle na validação por vaso.

As configurações dos três classificadores são apresentadas na
Tabela~\ref{tab:configuracao_modelos}.

\begin{table}[H]
    \centering
    \small
    \caption{Configurações dos modelos de aprendizado de máquina.}
    \label{tab:configuracao_modelos}
    \begin{tabular}{p{0.13\linewidth}p{0.22\linewidth}p{0.54\linewidth}}
        \hline
        \textbf{Modelo} & \textbf{Aspecto} & \textbf{Configuração} \\
        \hline
        ET & Implementação &
        \texttt{ExtraTreesClassifier}; 900 árvores;
        \texttt{max\_features="sqrt"}; mínimo de dez amostras por folha;
        \texttt{bootstrap=False}; sem pesos de classe; semente 42. \\
        \hline
        RF & Implementação &
        \texttt{RandomForestClassifier}; 300 árvores;
        \texttt{class\_weight="balanced"}; semente 42. Utilizado na
        comparação LOSO após as exclusões. \\
        \hline
        TCN & Arquitetura &
        Três blocos convolucionais temporais residuais, com 24, 24 e
        32 canais; duas convoluções de núcleo 3 por bloco; dilatações
        1, 2 e 4; agregação média global no tempo; \emph{dropout} de
        0,20; camada linear de duas classes. \\
        & Treinamento &
        Janelas de 128 leituras; 60 épocas; lotes de 32 janelas;
        otimizador AdamW; taxa de aprendizado $10^{-3}$;
        \emph{weight decay} de $10^{-4}$; entropia cruzada com pesos
        de valor inverso à frequência das classes nas
        janelas de treinamento; semente 42. \\
        & Predição por leitura &
        Média das probabilidades das janelas que cobriam cada leitura,
        com limiar de 0,5 para a classe 1. \\
        \hline
    \end{tabular}

    \vspace{0.1cm}
    {\footnotesize Fonte: Elaborado pelos autores (2026).}
\end{table}

Na TCN, as duas convoluções em cada bloco resultam em um campo
receptivo local de $1+4(1+2+4)=29$ leituras, equivalente a cerca de
6,7~s para o intervalo mediano de 0,24~s entre leituras. A agregação
média global combina as saídas locais de toda a janela de 128 leituras.
As 60 épocas e os demais hiperparâmetros foram mantidos fixos, sem
seleção por validação interna ou parada antecipada.

Na análise histórica, ET e TCN foram ajustados uma vez com as coletas
de treinamento. Solo seco e solo úmido foram dois recortes do mesmo
conjunto de teste, não dois treinamentos independentes. Nas análises
LOSO, os modelos foram reajustados em cada rodada, conforme descrito
adiante.

\paragraph{Análise histórica com separação por coleta}
\label{subsec:treinamento_teste}

A divisão foi feita por \emph{coleta}, dentro de cada classe, com
proporção nominal de 70\% para treinamento e 30\% para teste e semente 42.
Todas as leituras e janelas de uma coleta permaneceram em um único conjunto.
Das 36 coletas, 24 foram destinadas ao treinamento (41.377 leituras: 11
coletas da classe 0 e 13 da classe 1) e 12 ao teste (22.115 leituras:
seis coletas de cada classe).

Essa separação evita que janelas da mesma aquisição apareçam nos dois
conjuntos. Entretanto, oito dos dez vasos aparecem em ambos porque foram
medidos em dias distintos. Assim, o teste avalia novas \emph{coletas}, não vasos
independentes do ponto de vista biológico, e as métricas não demonstram generalização
para plantas ou vasos nunca vistos pelo modelo. Essa partição foi mantida
como comparação histórica; a validação agrupada por vaso foi usada como
referência principal para as análises com sinais MQ brutos.

No teste, o recorte seco incluiu 16.692 leituras de oito coletas: 6.507
leituras da classe 0 e 10.185 da classe 1. O recorte úmido incluiu 5.423
leituras de quatro coletas: 4.152 da classe 0 e 1.271 da classe 1. Essa
distribuição desigual motivou o relato da acurácia balanceada junto da
acurácia convencional. Como análise adicional da partição histórica,
foram usadas previsões já obtidas no teste para formar um subconjunto
com igual número de leituras das duas classes em cada faixa: 6.507 de
cada classe na faixa seca e 1.271 de cada classe na faixa úmida. Todas
as leituras da classe minoritária foram mantidas e a classe majoritária
foi subamostrada com semente 42; as leituras excedentes foram examinadas
em análise própria. Essa subamostragem ocorreu após a predição,
sem retreinamento, e foi tratada como análise de sensibilidade a um
sorteio específico, não como resultado principal.

\paragraph{Métricas da análise histórica}
\label{subsec:metricas}

Na partição histórica foram calculadas acurácia, acurácia balanceada,
F1 macro e matrizes de confusão no conjunto de teste. A acurácia é a
fração de leituras
classificadas com acerto; a acurácia balanceada é a média dos recalls
das classes 0 e 1; e a F1 macro é a média simples das medidas F1 das duas
classes. O recall da classe inoculada corresponde à sensibilidade
para a classe 0, ou seja, à fração de leituras de vasos inoculados
identificadas com acerto. As métricas foram calculadas para os recortes seco e
úmido do mesmo teste. A avaliação por subamostragem balanceada foi relatada
fora das métricas do teste completo.

\paragraph{Sequência das análises}
\label{subsec:sequencia_analises}

A Tabela~\ref{tab:cronologia_analises} distingue as bases e as decisões
tomadas em cada etapa. A busca de subconjuntos e a inspeção dos erros
ocorreram antes da análise pós-exclusão; os controles de entradas fixas
foram acrescentados depois para examinar a dependência dos
sinais ambientais.

\begin{table}[H]
    \centering
    \footnotesize
    \caption{Sequência das análises, base utilizada e decisão correspondente.}
    \label{tab:cronologia_analises}
    \begin{tabular}{p{0.07\textwidth}p{0.20\textwidth}p{0.61\textwidth}}
        \hline
        Etapa & Base & Objetivo e decisão \\
        \hline
        1 & 36 coletas & Recorte temporal e filtro estrito de pressão; definição da base comum. \\
        2 & 36 coletas & Partição histórica por coleta; comparação dos modelos com MQ corrigidos e avaliação adicional por condição do solo. \\
        3 & 36 coletas & Busca LOSO de 63 subconjuntos MQ com solo, temperatura e pressão fixos; seleção do conjunto MQ2, MQ3, MQ8, MQ135 e MQ138 e obtenção das previsões fora de cada vaso. \\
        4 & 36 coletas & Auditoria de saltos e de erros por coleta; definição exploratória das cinco coletas excluídas apenas da análise de sensibilidade. \\
        5 & 31 coletas & LOSO após exclusões; comparação de ET, Random Forest e TCN em entradas fixas e busca ET das 511 combinações. \\
        6 & 36 coletas & Controles LOSO com todas as nove entradas, MQ sem variáveis ambientais, variáveis ambientais sem MQ e regressão logística, sem seleção de subconjunto nessa etapa. \\
        7 & 36 coletas ou recortes por dia/tempo & Sensibilidade por dia, período e fase temporal; comparação pareada das configurações fixas. \\
        \hline
    \end{tabular}
\end{table}


\paragraph{Auditoria e exclusões exploratórias}
\label{subsec:auditoria_coletas_loso}

Na base completa, foram examinados os saltos entre leituras consecutivas dos
seis sensores MQ. Para cada sensor, o limiar de salto foi o percentil 99,5 das
diferenças absolutas calculadas em todas as coletas. Uma leitura foi marcada
quando pelo menos dois sensores superaram seus limiares. Uma coleta foi
sinalizada pela auditoria quando pelo menos 5\% de suas leituras receberam
essa marca. Os limiares foram definidos com a base completa, incluindo as
coletas reservadas para teste nas rodadas LOSO.

Na busca com as 36 coletas, o ET avaliou os 63 subconjuntos não vazios dos
seis MQ, mantendo solo, temperatura e pressão em cada configuração. As
previsões fora do vaso de treinamento foram usadas para identificar quatro
coletas com mais de 50\% de erros por leitura. Na análise de sensibilidade,
essas quatro coletas e uma quinta coleta sinalizada pela auditoria
multissensor foram retiradas: controle, vaso 2, em 12 de maio; inoculado,
vasos 9 e 10, em 19 de maio; controle, vasos 1 e 2, em 19 de maio.
A coleta do vaso 10 foi selecionada pela auditoria; as outras quatro foram
selecionadas pela fração de erros. A coleta controle do vaso 1 em 19 de maio
também apresentou saltos multissensor. A seleção ocorreu após a inspeção das
previsões e não altera a validação principal com 36 coletas. Saltos nos
sinais não confirmam falha instrumental.

\paragraph{Validação principal com coletas agrupadas por vaso}
\label{subsec:loso_complementar}

Para avaliar coletas de vasos não vistos durante o ajuste, aplicou-se a
validação \textit{leave-one-subject-out} (LOSO) à base completa,
com cada vaso como unidade de agrupamento e sem excluir coletas. Todas as
coletas de um mesmo vaso foram agrupadas:
em cada rodada, as coletas de um vaso constituíram o teste e as dos outros
nove constituíram o treinamento. Na base derivada, um vaso da classe
inoculada aparece com código numérico 1, embora a planilha de origem o
identifique como vaso 10. Por isso, a chave técnica de agrupamento combinou
classe e código de vaso, evitando colisão com o vaso 1 controle e mantendo
dez grupos distintos.
A base completa e a base da análise de sensibilidade
foram avaliadas em análises distintas com essa mesma regra. No segundo cenário,
retiraram-se apenas as cinco coletas descritas acima, mantendo os dez
vasos e suas demais coletas. Em nenhuma rodada coletas do vaso de teste
entraram no treinamento; cada leitura incluída foi testada uma vez.

Na base completa, o ET foi avaliado com MQ2, MQ3, MQ8, MQ135 e MQ138 mais
o índice de solo, a temperatura e a pressão. Essa configuração havia sido
escolhida entre subconjuntos de MQ avaliados nos mesmos vasos; o LOSO da
base completa é uma avaliação interna, não uma validação externa da
escolha das entradas.
O ET foi reajustado em
cada rodada com 900 árvores,
\texttt{max\_features="sqrt"}, \texttt{min\_samples\_leaf=10},
\texttt{bootstrap=False} e semente 42. A previsão por leitura foi a classe
com probabilidade estimada de classe 1 igual ou superior a 0,5.

Para distinguir a contribuição das entradas sem seleção de
um subconjunto, o mesmo LOSO na base completa foi repetido com três
configurações fixas de ET: os seis MQ mais solo, temperatura e pressão;
os seis MQ sem variáveis ambientais; e solo, temperatura e pressão sem MQ. Também se
avaliou, com as nove entradas, uma regressão logística como modelo
simples de referência. Nessa regressão, a padronização das entradas foi
ajustada nas coletas de treinamento de cada rodada; usaram-se
\texttt{LogisticRegression}, \texttt{max\_iter=2000}, sem pesos de classe,
e semente 42. Esses controles compartilham as mesmas dez divisões por
vaso e os mesmos pré-processamentos globais já descritos; não constituem
uma validação externa do protocolo experimental.

Na base após exclusão das cinco coletas, ET, \textit{Random Forest}
(\texttt{RandomForestClassifier}) e TCN foram reajustados em cada rodada
LOSO em duas configurações fixas: MQ2, MQ3, MQ8, MQ135 e MQ138 sem variáveis
ambientais; e os mesmos cinco MQ com solo, temperatura e pressão. O Random Forest usou 300
árvores, \texttt{class\_weight="balanced"} e semente 42. Na rede temporal,
a padronização das entradas foi estimada apenas no treinamento de cada
rodada. As janelas mantiveram 128 leituras e passo de 64, com cinco ou oito
canais conforme a configuração; nenhuma janela atravessou o limite de
uma coleta.
Os modelos foram comparados por leitura após reunir as previsões de todas
as rodadas. Essa avaliação difere da partição histórica de 70\%/30\% por coleta.
No treinamento e no teste LOSO, não houve
subamostragem, sobreamostragem nem geração de leituras sintéticas para
igualar o tamanho das classes. A subamostragem da partição histórica
ocorreu apenas depois das previsões. O ET
foi treinado sem pesos de classe; o Random Forest usou
\texttt{class\_weight="balanced"}, e a TCN usou pesos com valor inverso
à frequência das classes nas janelas de treinamento.
Portanto, os três modelos não receberam o mesmo tratamento do
desbalanceamento. A acurácia balanceada, o F1 macro e os recalls por
classe foram relatados junto da acurácia simples para avaliar as duas
classes em separado; essas métricas não modificam o treinamento do ET.

\paragraph{Busca combinatória das entradas}
\label{subsec:busca_combinatoria_loso}

A busca exaustiva de entradas foi realizada com ET, na base após
as cinco exclusões, usando a mesma divisão LOSO por vaso. Ela não foi
aplicada ao Random Forest nem à TCN.
Solo, temperatura e pressão variaram como entradas independentes nessa busca,
inclusive nas configurações sem sensores MQ. O BMP280 fornece
temperatura e pressão como duas variáveis, embora seja um único componente
físico; portanto, as 511 configurações correspondem a combinações de
variáveis de entrada, não de componentes.


As combinações foram ordenadas pela acurácia por leitura. Em caso de
empate, preferiu-se a configuração com menos variáveis; persistindo
o empate, a que continha mais sensores MQ e, em seguida, a ordem
alfabética dos nomes das entradas. Como análise complementar do mesmo
ranking, a acurácia balanceada de cada combinação foi calculada a partir
dos recalls das duas classes
sem novo treinamento dos modelos. A escolha original da configuração
utilizou a acurácia por leitura; o ranking balanceado foi examinado como
análise de sensibilidade. A escolha da configuração
e sua avaliação ocorreram nos mesmos vasos, sem validação aninhada
nem conjunto biológico externo; portanto, a ordenação das
configurações é exploratória.

\paragraph{Métricas da validação por vaso}
\label{subsec:metricas_loso_complementar}

O recall da classe 0 descreve a identificação de leituras do grupo
inoculado; nas métricas por leitura, os denominadores são leituras, não
vasos. As leituras sucessivas de uma mesma coleta e as coletas repetidas do mesmo
vaso não constituem réplicas biológicas independentes. Em particular,
a acurácia reunida atribui mais peso a coletas com maior número de leituras.

Para obter uma probabilidade por coleta, calculou-se a média das
probabilidades das suas leituras. Para obter uma probabilidade por vaso,
calculou-se a média das probabilidades de suas coletas, com o mesmo peso para
cada coleta. As decisões de coleta e vaso utilizaram limiar de 0,5 para a
classe 1.
As métricas foram calculadas nas três unidades:
leitura, coleta e vaso. Na unidade vaso, a proporção de acertos foi
expressa como $k/10$ e acompanhada por intervalo binomial exato de
95\% pelo método de Clopper--Pearson. Como medida exploratória adicional,
para a acurácia balanceada foi calculado um intervalo por 10.000
reamostragens estratificadas com reposição de
cinco vasos por classe; cada vaso levou consigo todas as suas previsões
fora de treino. Usaram-se os percentis 2,5 e 97,5 das estimativas
reamostradas, com semente 42. Esses intervalos são condicionais ao modelo
e às entradas já escolhidas, sem incorporar a incerteza da seleção da
configuração entre os subconjuntos avaliados. Com apenas dez vasos, o
intervalo por reamostragem foi tratado como descritivo.

Para os controles fixos de ET e a regressão logística na base completa,
assim como para a comparação ET, Random Forest e TCN na base com cinco
coletas excluídas, calculou-se a área sob a curva ROC (AUC) por coleta a
partir das probabilidades médias de cada coleta. Essa métrica não depende do
limiar de classificação de 0,5. Como as probabilidades de cada vaso
provêm de um modelo ajustado em uma rodada diferente, a AUC reunida
entre rodadas foi tratada como descritiva. Não foi realizado teste pareado
por vaso para afirmar superioridade entre ET, Random Forest e TCN.

\paragraph{Análises de sensibilidade temporal}
\label{subsec:sensibilidade_dia_periodo}

O ET foi reajustado com as 36 coletas em três separações: LOSO restrito aos
dias 18 e 19 de maio; treino nos dias 11 a 13 e teste nos dias 18 e 19; e
treino nos dias 18 e 19 e teste nos dias 11 a 13. Foram avaliadas três
configurações fixas de entradas: os seis MQ com solo, temperatura e pressão;
os seis MQ sem variáveis ambientais; e solo, temperatura e pressão sem MQ.
As separações entre períodos incluem vasos presentes no treino e no teste;
elas avaliam transferência temporal, não vasos novos.

Para avaliar a deriva entre dias, repetiu-se o LOSO com os seis MQ após
padronização por dia. Em cada rodada, a média e o desvio padrão de cada
sensor foram estimados nas leituras dos nove vasos de treino daquele dia e
aplicados às leituras de treino e teste. Outra análise LOSO reteve as
leituras com tempo decorrido entre 200 e 800~s após o filtro de pressão.
Esse intervalo fixo não identifica a exposição ao \textit{headspace} em todas
as coletas, pois as séries não contêm uma marca verificada da abertura da
válvula. As 36 coletas e os dez vasos permaneceram nessas análises.

Na base completa, as configurações de nove entradas e de três variáveis
ambientais foram comparadas por diferenças de acurácia em cada vaso. A
enumeração das $2^{10}$ inversões de sinal dessas diferenças foi usada como
análise pareada exploratória. Não houve permutação dos rótulos dos vasos com
novo treinamento dos modelos.
